\documentclass[11pt,a4paper]{article}
\usepackage[utf8]{inputenc}
\usepackage[T1]{fontenc}
\usepackage[english]{babel}
\usepackage{amsmath,amssymb}
\usepackage[round]{natbib}
\usepackage{booktabs}
\usepackage{graphicx}
\usepackage[margin=2.5cm]{geometry}
\usepackage{hyperref}
\usepackage{xcolor}
\hypersetup{colorlinks=true,linkcolor=blue!50!black,citecolor=blue!50!black,urlcolor=blue!50!black}
\usepackage{caption}
\captionsetup{font=small,labelfont=bf}

\title{\bfseries Similarity Finds the Fact, Not the Version\\
\large A Co-Trained Order Stamp in the Key Resolves Version Conflict\\
in Persistent Model-Internal Memory}

\author{Maximiliano Speranza\\
\small Independent Researcher, Buenos Aires, Argentina\\
\small \href{https://orcid.org/0009-0005-0413-8554}{ORCID 0009-0005-0413-8554}\\
\small \texttt{maximiliano.speranza@gmail.com}}

\date{September 2026}

\begin{document}
\maketitle

\begin{abstract}
\noindent
A persistent archive co-trained inside a language model learns to retrieve the right fact and then
fails to know which \emph{version} of it currently holds. We measure that failure exactly: with one
version archived the model answers $0.9974$, and with two versions of the same key competing it drops
to $0.4576$, which is chance between the old value and the new one. The retrieval is intact and the
ordering is absent.

We show that the missing information fits in a small co-trained \textbf{order stamp added to the
archive key}, and that it is the ordering, not the extra capacity, that does the work. Across five
seeds, revised keys go from $0.4570 \pm 0.0289$ without a stamp to $\mathbf{0.9956 \pm 0.0029}$ with
one. The falsification control, an identically parameterised stamp whose value bears no relation to
the real write turn, gives $0.4768 \pm 0.0389$, within one standard deviation of the baseline. The
fifteen runs order perfectly by condition with no overlap between cells.

Three further results sharpen the claim. First, with three versions per key, each written in its own
sequence so that no recency trace survives in the content, asking for the \emph{superseded} value goes
from $0.2917$, which is chance among three, to $0.8316$. Second, when write turns are sampled at random
so that no fixed slot-to-role table can solve the task, the stamp still yields $0.9644$ on the current
version, which rules out the table shortcut. Third, and practically the most useful, a \emph{corrupted}
stamp is worse than no stamp at all, degrading even item identification from $0.9974$ to $0.8802$. A
wrong timestamp is not missing information, it is damaging information.

We also separate two capabilities that are usually conflated. Preferring the latest value and using
the order are learned at different times, the first saturating around $3{,}000$ steps and the second
only emerging near $4{,}000$ and saturating near $6{,}000$. A benchmark that only asks for the current
value cannot see the distinction, because the recency shortcut already answers it.
\end{abstract}

\section{The problem, measured}

Linear-attention and delta-rule memories have a fixed-size recurrent state and overwrite on write, so
what is old is lost by construction \citep{deltanet,gateddeltanet}. The natural repair is an external
archive, and a growing line of work co-trains that archive inside the model rather than bolting a
frozen retriever onto it \citep{colmlm,persistmem,lmlm}.

That repair works for retrieval and leaves a specific hole. In our own prior stage, a co-trained
archive holding vectors produced by the model itself reaches $0.9974$ when a key has a single archived
version, and collapses to $0.4576$ when the same key has two versions competing. The value $0.4576$ is
not noise, it is chance between the old and the new. The model finds the right fact and does not know
which version rules.

This is the same failure mode we had already measured in a non-parametric index, where geometry
grouped versions of a memory perfectly and recovered the \emph{oldest} one. Two mechanisms with
nothing in common fail in the same way, which suggests the diagnosis is a property of
similarity-based recall and not of any particular implementation. Similarity identifies the item.
Order is information similarity does not contain.

\section{Related work}

\textbf{Intra-sequence hybrids.} \citet{hola} pairs a linear-attention state with an exact memory for
what that state forgets, and \citet{ham} combines a recurrent state with attention that supplements it
only where the state predicts poorly. Both are hybrids of a compressive state with an exact store, and
in both the store lives inside one forward pass, so the question of which of two versions written in
\emph{different} sessions holds cannot arise.

\textbf{Co-trained persistent memory.} Co-LMLM \citep{colmlm} is the closest in mechanism, emitting
continuous queries from the model's hidden states to retrieve from a knowledge base, at $135$M to
$360$M parameters. Trained Persistent Memory \citep{persistmem} accumulates a memory bank across turns
on a frozen GPT-2 backbone so facts stated in one session are recalled later. Neither versions a fact,
neither stamps order into the key, and neither ever asks for a superseded value. Co-LMLM treats the
base as a static snapshot; \citet{persistmem} applies a uniform temporal decay factor, which
attenuates old entries rather than ordering them.

\textbf{Proactive interference.} \citet{unabletoforget} is the closest in \emph{task}, updating keys
several times within a context and asking for the current value with old versions present as
interference. It documents the same recency shortcut we find. It is a benchmark and a diagnosis, its
authors state that they offer no architectural mitigation, it evaluates only large pretrained models,
and it never asks for the previous value. That last point matters for us, because the recency shortcut
answers the current-value question by itself.

\textbf{Application-layer versioning.} Systems that version documents with explicit metadata and
pointers do exist, in retrieval pipelines built on top of a frozen model. The stamp we study is not
that: it lives inside the model's own retrieval key and is trained jointly with the reader, rather than
sitting in a metadata field consumed by a separate retriever.

\textbf{Timestamps in the input.} \citet{tslm} injects time markers into the input for event streams,
past, present and future. That is not two versions of the same fact inside a memory, and the marker
does not enter a retrieval key.

\section{Setup}

The substrate is a small language model with a co-trained archive, $64$ hidden dimensions, an archive
of $9$ entries, and a synthetic cross-sequence task. Facts are written in one sequence, revised in
another, and queried in a third, \textbf{with the recurrent state reset between sequences}. Retrieval
therefore cannot be solved by carrying activations forward, and the floor is chance by construction.

The archive stores vectors produced by the model, not tokens. Writing entry $i$ from state $h_i$ gives
key $k_i = h_i W_k$ and value $v_i = h_i W_v$; reading forms a query from the consuming state and
attends over the archive. The \textbf{stamp} adds a learned embedding of the write turn to the key,
\[
k_i \;=\; h_i W_k \;+\; \mathrm{ord}[\,t_i\,],
\]
where $t_i$ is the turn at which entry $i$ was written and $\mathrm{ord}$ is a trained table. It is
$384$ effective parameters. Crucially, \textbf{the archive is shuffled per sample} with each stamp
travelling attached to its own entry. Without shuffling, write turn would coincide with tensor
position and the experiment would measure ``look at the end of the archive'', which is position and
not order.

Three conditions are compared throughout. \texttt{none} has no stamp. \texttt{stamp} is as above.
\texttt{shuffled} is the falsification control, with an identically sized and identically trained
table whose stamp value bears \emph{no relation} to the real write turn.

\section{Result 1: the stamp lifts version conflict, and it is the ordering that lifts it}

Five seeds, two versions per revised key, $4{,}000$ steps.

\begin{center}
\begin{tabular}{lccc}
\toprule
condition & global & \textbf{revised keys} & unrevised keys \\
\midrule
\texttt{none} & $0.7273 \pm 0.0146$ & $0.4570 \pm 0.0289$ & $0.9977$ \\
\textbf{\texttt{stamp}} & $\mathbf{0.9970 \pm 0.0020}$ & $\mathbf{0.9956 \pm 0.0029}$ & $0.9984$ \\
\texttt{shuffled} & $0.7374 \pm 0.0194$ & $0.4768 \pm 0.0389$ & $0.9979$ \\
\bottomrule
\end{tabular}
\end{center}

The gain over no stamp is $+0.5385$ and over the shuffled control $+0.5187$. \textbf{The five stamped
seeds fall between $0.9909$ and $0.9987$, and the best seed of the other two conditions reaches
$0.5156$. The fifteen runs order perfectly by condition with no overlap.} At that separation a
confidence interval is a formality.

\textbf{The cell that did not win is what makes the result valid.} \texttt{shuffled} has exactly the
same parameters as \texttt{stamp}, trained identically, and differs only in that the stamp carries no
relation to the real write turn. Without it, $+0.5385$ would be explained just as well by saying we
gave the reader more capacity. With it, the capacity explanation is ruled out directly.

\textbf{The arithmetic is symmetric.} Version conflict cost $-0.5398$ and the stamp returns $+0.5385$.
It restores what similarity lost and nothing more, leaving unrevised keys at $0.998$ where they
already were.

\section{Result 2: three versions, and asking for the superseded value}

A first attempt at three versions broke itself, and the way it broke is worth recording. Writing both
revisions inside the same sequence let the state at $v_3$ carry a trace of $v_2$, so order was written
into the \emph{content} without any metadata putting it there. At $4{,}000$ steps this was invisible
because no condition solved the hard question; at $12{,}000$ steps the no-stamp condition reached
$0.97$ and exposed the leak.

With each version written in its own sequence and the state reset between them, $v_2$ and $v_3$ are
each ``first and only mention of their sequence'', indistinguishable in recency. The only thing in the
system that says which came first is the stamp. Three seeds, $12{,}000$ steps.

\begin{center}
\begin{tabular}{lccc}
\toprule
condition & current & \textbf{previous} & single-version keys \\
\midrule
\texttt{none} & $0.3090$ & $0.2917 \pm 0.0207$ & $0.9436$ \\
\textbf{\texttt{stamp}} & $0.9774$ & $\mathbf{0.8316 \pm 0.2399}$ & $0.9974$ \\
\texttt{shuffled} & $0.2682$ & $0.2804 \pm 0.0159$ & $0.8802$ \\
\bottomrule
\end{tabular}
\end{center}

Without a stamp the model sits at $1/3$, chance among three versions, while still identifying the item
when there is no conflict ($0.9436$). The leak is closed, and the distance from the $0.97$ of the
broken design measures exactly how much order had been leaking through content.

\textbf{This number is budget-dependent and we report it as such.} With five seeds the previous-value
score falls to $0.7667$ (sd $0.2340$): $0.9766$, $0.9635$, $0.5547$, $0.8594$, $0.4792$. The mechanism
does not fail, two of five seeds have not converged at $12{,}000$ steps. Tested directly, seed 2
extended to $24{,}000$ steps moved from $0.5547$ to $0.9531$. The honest verdict is conditional: the
stamp allows answering about the previous version, at a budget that $12{,}000$ steps does not always
reach. \textbf{In this regime convergence time is part of the phenomenon, not an implementation
detail.}

\section{Result 3: it is not a slot-to-role table}

In the previous design the write turns were fixed per role, so the model could learn ``embeddings
$9,10,11$ mean current''. That is a slot-to-role correspondence, not an order, and the shuffled control
does not separate them because randomising the stamp breaks the table just as it breaks the order.

We therefore sampled twelve turns at random out of a range of thirty-two per sample, sorted them, and
assigned them respecting the real write order. Turn $9$ can be a first version in one sample and a
third in the next, so no fixed table solves the task. Three seeds.

\begin{center}
\begin{tabular}{lcc}
\toprule
condition & \textbf{current} & previous \\
\midrule
\texttt{none} & $0.3247$ & $0.2977 \pm 0.0326$ \\
\textbf{\texttt{stamp}} & $\mathbf{0.9644}$ & $0.3142 \pm 0.5330$ \\
\texttt{shuffled} & $0.2804$ & $0.2925 \pm 0.0266$ \\
\bottomrule
\end{tabular}
\end{center}

The current-version result survives the moving turns, which \textbf{rules out the table shortcut and
validates Result 1 backwards}: had the mechanism been a table, moving the turns would have broken the
current version too.

The previous-value column is \textbf{bimodal, not noisy}: $0.0052$, $0.9297$, $0.0078$. Two seeds fall
into a recency shortcut, answering \emph{below} chance because they return the current value when asked
for the previous one, and one seed learns the full comparison. We report the distribution rather than
the mean, because the mean of $0.31$ describes none of the three runs. The honest summary is that the
reader learns to read the clock to know what time it is now, not to reconstruct the sequence of what
happened.

\section{Result 4: a lying stamp is worse than no stamp}

Reading the \texttt{shuffled} row of Result 2 as a result in itself, rather than only as a control,
gives the finding with the most practical bite. A corrupted stamp leaves the current version at
$0.2682$, \emph{below} the $0.3090$ of having no stamp at all, and it degrades item identification from
$0.9974$ to $0.8802$, a capability that has nothing to do with ordering.

\textbf{The reader leans on the stamp.} Once it does, a corrupted timestamp is not missing information,
it is damaging information, and the damage spreads to retrieval quality. For any system that stamps
memory entries with a write time, this says that a stale or wrong stamp is worse than an absent field.

\section{Two capabilities, learned at different times}

Preferring the latest value and using the order are distinct and separable, and the ordering of the
two is stable across replicates even though the timing is not. In both replicates at $lr=10^{-3}$ the
current version saturates by $3{,}000$ steps ($0.9583$ in both) while the previous-value score is
still $0.0000$. The previous-value score then lifts off later, and \emph{how much} later varies: one
replicate reaches $0.2500$ at $4{,}000$, $0.6250$ at $5{,}000$ and $0.9583$ at $6{,}000$, while the
other is still at $0.0208$ at $4{,}000$ and does not pass $0.9$ until $8{,}000$. We report both because
the mean of the two would describe neither. What replicates is the \textbf{ordering}: in every run the
current version is solved thousands of steps before the previous one, and at the step where the current
version is already at $0.9583$ the previous one is at exactly $0.0000$.

This has a direct consequence for evaluation. \textbf{A benchmark that only asks for the current value
cannot see the distinction}, because the recency shortcut gives the right answer. This is precisely
what \citet{unabletoforget} does not do, and it is why we consider the previous-value question the more
informative probe.

\section{Limitations}

\textbf{Sufficiency, not inference.} In this task the current version is always the one with the
highest turn, so the learnable policy is ``take the largest stamp''. Our result says the information
similarity lacks \emph{fits in a stamp} and that the reader can use it. It does not say the model
reasons about time.

\textbf{The stamp is given, not inferred.} We show the reader can use the metadata, not that it can
fabricate it. This is not an unrealistic oracle, since write turn is free to any index, unlike knowing
which version is correct, but it does bound the claim.

\textbf{Scale.} A $64$-dimensional model, a $9$-entry archive and a synthetic task. The measured
difficulty is alignment and version conflict, not discrimination among many candidates.

\textbf{A fresh archive.} The archive is recomputed at every step in these experiments. A follow-up in
which entries age shows the damage is gradual rather than a cliff, with revised keys falling from
$0.9987$ to $0.9115$ as the cosine between write-time and read-time frames falls to $0.78$.

\section{Conclusion}

Similarity identifies the item; order is information that similarity does not contain and that has to
be carried separately. In a persistent memory co-trained inside a small language model, carrying it as
a stamp in the key takes version conflict from chance to near ceiling, and a control with the same
parameters and a meaningless stamp shows the ordering rather than the capacity is doing the work. The
same mechanism makes the superseded value answerable, which no benchmark in this area currently asks
for, and it comes with a warning that generalises past our setup: once a reader learns to lean on a
stamp, a wrong stamp costs more than no stamp.

\section*{Reproducibility and data availability}

Every experiment reported here was pre-registered with its predictions and thresholds hashed before
the corresponding run. Cells that failed their pre-registered criteria are reported as failures,
including the budget-dependent previous-value score of Result 2 and the bimodal column of Result 3.
All pre-registrations with their SHA hashes, analysis scripts, per-seed raw results and
machine-generated reports are archived in the project repository,
\url{https://github.com/SperanzaMax/delta-archivo}, whose commit history carries the timestamps that
establish the order of pre-registration and result.

\section*{Declarations}

\textbf{Funding.} None. \textbf{Competing interests.} None. \textbf{Ethics approval.} Not applicable;
no human or animal subjects. \textbf{Data availability.} See above. \textbf{Author contributions.}
Sole author.

\begin{thebibliography}{99}
\bibitem[Cui(2026)]{hola} Cui, W. (2026). A Hippocampus for Linear Attention: An Exact Memory for What the Recurrent State Forgets. arXiv:2607.02303.
\bibitem[Lufkin et al.(2026)]{ham} Lufkin, L., Figliolia, T., Millidge, B., and Krishnamurthy, K. (2026). Hybrid Associative Memories. arXiv:2603.22325.
\bibitem[Feldman et al.(2026)]{colmlm} Feldman, Y., Zhao, L., Godey, N., Go, D., Hua, Y., Weinberger, K. Q., Sun, J. J., and Artzi, Y. (2026). Co-LMLM: Continuous-Query Limited Memory Language Models. arXiv:2607.07707.
\bibitem[Jeong(2026)]{persistmem} Jeong, H. (2026). Trained Persistent Memory for Frozen Decoder-Only LLMs. arXiv:2603.22329.
\bibitem[Zhao et al.(2025)]{lmlm} Zhao, L., Zalouk, S., Belardi, C. K., Lovelace, J., Zhou, J. P., Noonan, R. T., Go, D., Weinberger, K. Q., Artzi, Y., and Sun, J. J. (2025). Pre-training Limited Memory Language Models with Internal and External Knowledge. arXiv:2505.15962.
\bibitem[Wang and Sun(2025)]{unabletoforget} Wang, C., and Sun, J. V. (2025). Unable to Forget: Proactive Interference Reveals Working Memory Limits in LLMs Beyond Context Length. arXiv:2506.08184.
\bibitem[Schlag et al.(2021)]{deltanet} Schlag, I., Irie, K., and Schmidhuber, J. (2021). Linear Transformers Are Secretly Fast Weight Programmers. ICML 2021.
\bibitem[Yang et al.(2025)]{gateddeltanet} Yang, S., Kautz, J., and Hatamizadeh, A. (2025). Gated Delta Networks: Improving Mamba2 with Delta Rule. ICLR 2025. arXiv:2412.06464.
\bibitem[Rajaby Faghihi and Kordjamshidi(2021)]{tslm} Rajaby Faghihi, H., and Kordjamshidi, P. (2021). Time-Stamped Language Model: Teaching Language Models to Understand the Flow of Events. NAACL 2021. arXiv:2104.07635.
\end{thebibliography}

\end{document}
